% ============================================================================
\chapter{Tsetlin Machine Clause Learning}
\label{ch:clauses}
% ============================================================================

This is the core chapter of the Tsetlin Machine. It explains how
clauses are organised into a classifier, how feedback drives learning,
and how the entire system learns to recognise patterns. Read this
chapter carefully---the mechanisms here are the foundation for everything
that follows.


% ----------------------------------------------------------------------------
\section{Architecture Overview}
\label{sec:tm-architecture}
% ----------------------------------------------------------------------------

A Tsetlin Machine~\cite{granmo2018tsetlin} for $K$-class classification (in our case $K = 10$
for digits 0--9) is organised as follows:

\begin{enumerate}
  \item Each class $k \in \{0, 1, \ldots, K{-}1\}$ has $m$ clauses.
    Half are designated \emph{positive} ($C_1^+, C_2^+, \ldots,
    C_{m/2}^+$) and half \emph{negative} ($C_1^-, C_2^-, \ldots,
    C_{m/2}^-$).
  \item Each clause is a conjunctive formula (an AND of selected
    literals), controlled by a team of Tsetlin Automata as described
    in \cref{ch:automaton}.
  \item The \concept{clause vote} for class $k$ is:
    \begin{equation}
      v_k = \sum_{j=1}^{m/2} C_j^+ - \sum_{j=1}^{m/2} C_j^-
      \label{eq:clause-vote}
    \end{equation}
    Positive clauses vote \emph{for} the class; negative clauses
    vote \emph{against}.
  \item The predicted class is the one with the highest vote:
    \begin{equation}
      \hat{y} = \arg\max_k \; v_k
      \label{eq:predict}
    \end{equation}
\end{enumerate}

\begin{figure}[htbp]
\centering
\begin{tikzpicture}[
  clause/.style={draw, rounded corners, minimum width=1.5cm,
    minimum height=0.6cm, font=\scriptsize\sffamily, inner sep=2pt},
  posclause/.style={clause, fill=bitblue!15},
  negclause/.style={clause, fill=bitorange!15},
  >=Stealth
]
  % Input
  \node[font=\small\sffamily\bfseries] at (-2, 0) {Input $\bvec{x}$};

  % Class 0 clauses
  \node[font=\small\sffamily\bfseries] at (3, 2.5) {Class 0};
  \node[posclause] (c0p1) at (1.5, 1.8) {$C_1^+$};
  \node[posclause] (c0p2) at (3.0, 1.8) {$C_2^+$};
  \node[negclause] (c0n1) at (4.5, 1.8) {$C_1^-$};
  \node[negclause] (c0n2) at (6.0, 1.8) {$C_2^-$};

  % Vote
  \node[draw, circle, minimum size=0.8cm, font=\small\sffamily,
    fill=bitgreen!10] (v0) at (8, 1.8) {$v_0$};

  \draw[->, thick] (c0p1) -- (v0) node[midway, above, font=\tiny] {$+$};
  \draw[->, thick] (c0p2) -- (v0);
  \draw[->, thick] (c0n1) -- (v0) node[midway, below, font=\tiny] {$-$};
  \draw[->, thick] (c0n2) -- (v0);

  % Class 1 (abbreviated)
  \node[font=\small\sffamily\bfseries] at (3, 0.3) {Class 1};
  \node[font=\scriptsize\sffamily, bitgray] at (3.75, -0.3) {$\cdots$ (same structure)};
  \node[draw, circle, minimum size=0.8cm, font=\small\sffamily,
    fill=bitgreen!10] (v1) at (8, -0.3) {$v_1$};

  % Dots
  \node[font=\sffamily] at (3, -1.2) {$\vdots$};
  \node[font=\sffamily] at (8, -1.2) {$\vdots$};

  % Class 9
  \node[font=\small\sffamily\bfseries] at (3, -2.1) {Class 9};
  \node[font=\scriptsize\sffamily, bitgray] at (3.75, -2.7) {$\cdots$};
  \node[draw, circle, minimum size=0.8cm, font=\small\sffamily,
    fill=bitgreen!10] (v9) at (8, -2.7) {$v_9$};

  % Argmax
  \node[draw, rounded corners, fill=bitpurple!10, minimum width=1.5cm,
    font=\small\sffamily\bfseries] (argmax) at (10.5, -0.3)
    {$\arg\max$};
  \draw[->, thick] (v0) -- (argmax);
  \draw[->, thick] (v1) -- (argmax);
  \draw[->, thick] (v9) -- (argmax);

  \node[font=\small\sffamily\bfseries, right] at (11.5, -0.3) {$\hat{y}$};
\end{tikzpicture}
\caption{Tsetlin Machine architecture for 10-class classification.
  Each class has $m$ clauses (half positive, half negative). The
  clause vote is the sum of positive clause outputs minus negative
  clause outputs. The predicted class is the argmax of all votes.}
\label{fig:tm-architecture}
\end{figure}


% ----------------------------------------------------------------------------
\section{Clause Evaluation}
\label{sec:clause-eval}
% ----------------------------------------------------------------------------

Given an input $\bvec{x} \in \{0, 1\}^n$ (after booleanisation and
literal expansion, so $n = 2 \times 784 = 1568$ for MNIST), a clause
with included literal set $I \subseteq \{1, \ldots, n\}$ evaluates as:

\begin{equation}
  C(\bvec{x}) = \bigwedge_{k \in I} \ell_k(\bvec{x})
  \label{eq:clause-eval}
\end{equation}

If $I = \emptyset$ (no literals included), the clause outputs 1 for
all inputs---it is vacuously true. The training algorithm prevents
this degenerate case through Type~I feedback (\cref{sec:type1}).

\begin{engineernote}
Clause evaluation maps directly to bitwise AND on packed bit vectors.
For MNIST's 1568 literals, we need $\lceil 1568 / 64 \rceil = 25$
64-bit words. The clause's ``include mask'' is also 25 words. The
evaluation is:
\begin{enumerate}
  \item Bitwise AND the include mask with the literal vector.
  \item Compare the result to the include mask.
  \item If they are equal (all included literals are 1), the clause
    outputs 1.
\end{enumerate}
This takes roughly 50 instructions regardless of how many literals
are included.
\end{engineernote}

In practice, the evaluation is often written as:
\begin{equation}
  C(\bvec{x}) = \begin{cases}
    1 & \text{if } (\bvec{m} \wedge \bvec{x}) = \bvec{m} \\
    0 & \text{otherwise}
  \end{cases}
  \label{eq:clause-bitmask}
\end{equation}
where $\bvec{m}$ is the clause's bitmask with 1s at positions of
included literals.


% ----------------------------------------------------------------------------
\section{The Voting Mechanism and Thresholding}
\label{sec:voting}
% ----------------------------------------------------------------------------

The raw clause vote $v_k$ from \cref{eq:clause-vote} is typically
\concept{clipped} to a threshold $T$:

\begin{equation}
  v_k^{\text{clipped}} = \max(-T, \; \min(T, \; v_k))
  \label{eq:clipped-vote}
\end{equation}

The threshold $T$ limits how strongly any single class can dominate.
It also plays a role in feedback: training feedback is only generated
when $|v_k| < T$, so clauses stop being reinforced once the class
vote is already strong enough. This prevents over-learning.

\begin{keyinsight}
The threshold $T$ controls the number of clauses that actively
participate in learning. If $v_k$ already exceeds $T$, additional
positive clauses are not reinforced---this frees up ``training
capacity'' for other patterns. Think of it as a budget: each class
can spend at most $T$ votes of confidence, encouraging diverse
clause patterns rather than redundant ones.
\end{keyinsight}


% ----------------------------------------------------------------------------
\section{The Feedback Mechanism}
\label{sec:feedback}
% ----------------------------------------------------------------------------

Training a Tsetlin Machine means adjusting the states of all the
automata so that clauses learn to detect useful patterns. The feedback
mechanism is the heart of the algorithm. It comes in two types.

Given a training example $(\bvec{x}, y)$ where $y$ is the true class:

\begin{description}[leftmargin=1em]
  \item[Type~I feedback] is given to clauses of the \textbf{correct
    class} $y$. Its purpose is to make these clauses \emph{more
    specific}---to include more literals that match the current input
    and exclude literals that do not.

  \item[Type~II feedback] is given to clauses of an
    \textbf{incorrect class} $k \neq y$ (specifically, the class
    with the highest or a randomly chosen high vote). Its purpose is
    to make these clauses \emph{less matching}---to push them away
    from recognising the current input.
\end{description}


% ----------------------------------------------------------------------------
\section{Type~I Feedback: Reinforcing the Correct Class}
\label{sec:type1}
% ----------------------------------------------------------------------------

Type~I feedback is given to clauses belonging to the true class $y$.
It operates differently on positive and negative clauses, and
differently depending on whether the clause matched the current input.

\subsection{Type~I on Positive Clauses ($C^+$)}

The goal is to strengthen positive clauses that match patterns from
class $y$.

\paragraph{Case 1: The clause output is 1 (clause matches the input).}

For each literal $\ell_k$ in the clause:
\begin{itemize}
  \item If $\ell_k = 1$ (the literal is true in this input) and the
    automaton says ``include'': \textbf{reward} the automaton with
    probability $\frac{s-1}{s}$. This reinforces the inclusion of
    matching literals.
  \item If $\ell_k = 0$ (the literal is false in this input) and the
    automaton says ``include'': this means the clause includes a
    literal that is currently 0, but the clause still output 1. This
    cannot actually happen (the clause is an AND of included
    literals, so all included literals must be 1 for the clause to
    output 1). So this case is impossible.
  \item If $\ell_k = 1$ and the automaton says ``exclude'':
    \textbf{penalty} the automaton with probability $\frac{s-1}{s}$.
    This pushes excluded-but-true literals toward inclusion,
    making the clause more specific.
  \item If $\ell_k = 0$ and the automaton says ``exclude'':
    \textbf{reward} the automaton with probability $\frac{1}{s}$.
    This gently reinforces the exclusion of currently-false literals.
\end{itemize}

\paragraph{Case 2: The clause output is 0 (clause does not match the input).}

For each literal:
\begin{itemize}
  \item If the automaton says ``include'': \textbf{penalty} with
    probability $\frac{1}{s}$. This gently pushes included literals
    toward exclusion, making the clause less specific (so it might
    match more inputs in the future).
  \item If the automaton says ``exclude'': \textbf{reward} with
    probability $\frac{1}{s}$. This gently reinforces exclusion.
\end{itemize}

\begin{engineernote}
Notice the asymmetry: when the clause matches ($C = 1$), feedback is
applied with high probability $\frac{s-1}{s}$. When the clause does
not match ($C = 0$), feedback uses the low probability $\frac{1}{s}$.
This means matching clauses are strongly reinforced while non-matching
clauses drift slowly. The specificity parameter $s$ controls this
ratio.
\end{engineernote}

\subsection{Type~I on Negative Clauses ($C^-$)}

Negative clauses vote \emph{against} the class. For the true class
$y$, we want negative clauses to \emph{not match}. So Type~I feedback
on negative clauses is the mirror image: it tries to push them away
from the current input pattern.

In practice, Type~I feedback on negative clauses for the true class
is either omitted or applies only the ``clause output is 0'' case
(gently reinforcing exclusion). The key training signal for negative
clauses comes from Type~II feedback on incorrect classes.


% ----------------------------------------------------------------------------
\section{Type~II Feedback: Weakening Incorrect Classes}
\label{sec:type2}
% ----------------------------------------------------------------------------

Type~II feedback targets clauses of an \textbf{incorrect} class
$k \neq y$ that have high votes for the current input. The goal is to
break these clauses so they stop matching.

\paragraph{When does Type~II apply?}
For a positive clause $C^+$ of class $k \neq y$ that outputs~1 (it
wrongly matches the input):

For each literal $\ell_k$ that is \textbf{included} in the clause:
\begin{itemize}
  \item If $\ell_k = 0$ (the literal is false in this input): this
    cannot happen (clause output is 1, so all included literals are
    true).
  \item If $\ell_k = 1$ (the literal is true): \textbf{penalty} the
    automaton, pushing it from ``include'' toward ``exclude.''
\end{itemize}

For each literal that is \textbf{excluded}:
\begin{itemize}
  \item If $\ell_k = 0$ (the literal is false): \textbf{reward} the
    automaton (reinforcing exclusion of this false literal---if it
    were included, it would break the clause, which is what we want,
    but we also want the automaton to ``remember'' this exclusion).
\end{itemize}

\begin{keyinsight}
Type~II feedback has a specific and elegant effect: it tries to
include literals that are 0 in the current input. If it succeeds in
including even one such literal, the clause will output 0 for this
input (because the AND of a 0 is 0), breaking the false match. This
is a \emph{targeted sabotage}: instead of destroying the entire
clause, Type~II just adds one ``poison pill'' literal that prevents
the clause from matching this particular wrong input.
\end{keyinsight}

\begin{figure}[htbp]
\centering
\begin{tikzpicture}[
  >=Stealth,
  box/.style={draw, rounded corners, minimum width=3cm, minimum height=1cm,
    font=\small\sffamily, align=center}
]
  % Type I
  \node[box, fill=bitblue!10] (t1) at (0, 0)
    {\textbf{Type~I Feedback}\\Correct class $y$};
  \node[font=\scriptsize\sffamily, text width=4cm, align=center]
    at (0, -1.5)
    {Strengthen matching clauses\\Include true literals\\Exclude false literals};

  % Type II
  \node[box, fill=bitorange!10] (t2) at (7, 0)
    {\textbf{Type~II Feedback}\\Incorrect class $k \neq y$};
  \node[font=\scriptsize\sffamily, text width=4cm, align=center]
    at (7, -1.5)
    {Break wrongly matching clauses\\Include a false literal\\to sabotage the clause};

  % Arrows
  \draw[->, thick, bitblue] (-2, 1.5) -- (t1)
    node[midway, left, font=\scriptsize\sffamily] {train example};
  \draw[->, thick, bitorange] (9, 1.5) -- (t2)
    node[midway, right, font=\scriptsize\sffamily] {train example};
\end{tikzpicture}
\caption{The two types of feedback. Type~I reinforces clauses of the
  correct class, making them more specific. Type~II sabotages clauses
  of incorrect classes that wrongly matched the input.}
\label{fig:feedback-types}
\end{figure}


% ----------------------------------------------------------------------------
\section{The Complete Feedback Tables}
\label{sec:feedback-tables}
% ----------------------------------------------------------------------------

For reference, we present the complete feedback rules in tabular form.
Let $C$ be the clause output, $\ell$ the literal value in the current
input, and $a$ the automaton action (0 = exclude, 1 = include).

\subsection{Type~I Feedback (Positive Clause, True Class)}

\begin{table}[htbp]
\centering
\caption{Type~I feedback for a positive clause of the true class.
  ``$\nearrow$'' means reward (push away from boundary);
  ``$\searrow$'' means penalty (push toward boundary);
  ``---'' means no feedback. Probabilities are shown in parentheses.}
\label{tab:type1-pos}
\begin{tabular}{@{}ccc|l@{}}
  \toprule
  $C$ & $\ell$ & $a$ & Feedback \\
  \midrule
  1 & 1 & 1 (include) & $\nearrow$ Reward with prob.\ $\frac{s-1}{s}$ \\
  1 & 1 & 0 (exclude) & $\searrow$ Penalty with prob.\ $\frac{s-1}{s}$ \\
  1 & 0 & 1 (include) & impossible ($C = 1$ requires all included $\ell = 1$) \\
  1 & 0 & 0 (exclude) & $\nearrow$ Reward with prob.\ $\frac{1}{s}$ \\
  \midrule
  0 & 1 & 1 (include) & $\searrow$ Penalty with prob.\ $\frac{1}{s}$ \\
  0 & 1 & 0 (exclude) & $\nearrow$ Reward with prob.\ $\frac{1}{s}$ \\
  0 & 0 & 1 (include) & $\searrow$ Penalty with prob.\ $\frac{1}{s}$ \\
  0 & 0 & 0 (exclude) & $\nearrow$ Reward with prob.\ $\frac{1}{s}$ \\
  \bottomrule
\end{tabular}
\end{table}

\subsection{Type~II Feedback (Positive Clause, Wrong Class)}

\begin{table}[htbp]
\centering
\caption{Type~II feedback for a positive clause of an incorrect class.
  Only clauses with $C = 1$ receive Type~II feedback.}
\label{tab:type2-pos}
\begin{tabular}{@{}ccc|l@{}}
  \toprule
  $C$ & $\ell$ & $a$ & Feedback \\
  \midrule
  1 & 1 & 1 (include) & --- (no action) \\
  1 & 1 & 0 (exclude) & $\searrow$ Penalty (push toward include) \\
  1 & 0 & 1 (include) & impossible \\
  1 & 0 & 0 (exclude) & $\searrow$ Penalty (push toward include) \\
  \midrule
  0 & \multicolumn{2}{c|}{any} & --- (no feedback) \\
  \bottomrule
\end{tabular}
\end{table}

\begin{engineernote}
The tables above describe the most common formulation
from~\cite{granmo2018tsetlin}. Variants exist: some
implementations apply Type~II feedback deterministically (probability
1), others use $\frac{1}{s}$. The variant matters less than the
principle: Type~I builds, Type~II sabotages.
\end{engineernote}


% ----------------------------------------------------------------------------
\section{Putting It All Together: One Training Step}
\label{sec:training-step}
% ----------------------------------------------------------------------------

\begin{algorithm}[htbp]
\caption{One training step of the Tsetlin Machine.}
\label{alg:tm-train}
\begin{algorithmic}[1]
  \State \textbf{Input:} Training example $(\bvec{x}, y)$, threshold $T$,
    specificity $s$
  \State Booleanise input: $\bvec{x}_{\text{bool}} \leftarrow
    \text{threshold\_and\_expand}(\bvec{x})$
  \For{each class $k = 0, 1, \ldots, K{-}1$}
    \State Evaluate all clauses: compute $C_j^+(\bvec{x}_{\text{bool}})$
      and $C_j^-(\bvec{x}_{\text{bool}})$
    \State Compute vote: $v_k \leftarrow \sum C_j^+ - \sum C_j^-$
    \State Clip: $v_k \leftarrow \text{clamp}(v_k, -T, T)$
  \EndFor
  \Statex
  \For{each class $k = 0, 1, \ldots, K{-}1$}
    \If{$k = y$} \Comment{True class}
      \State Apply Type~I feedback to clauses of class $k$
      \Statex \quad\quad with probability $\frac{T - v_k}{2T}$
    \Else \Comment{Wrong class}
      \State Apply Type~II feedback to clauses of class $k$
      \Statex \quad\quad with probability $\frac{T + v_k}{2T}$
    \EndIf
  \EndFor
\end{algorithmic}
\end{algorithm}

The probabilities on lines 10 and 13 deserve attention:

\begin{itemize}
  \item For the true class $y$: feedback probability is
    $\frac{T - v_y}{2T}$. When $v_y$ is low (the correct class is
    losing), this probability is high---urgent reinforcement. When
    $v_y \approx T$ (the correct class is winning strongly), the
    probability drops to near zero---no need to train further.

  \item For wrong classes $k \neq y$: feedback probability is
    $\frac{T + v_k}{2T}$. When $v_k$ is high (the wrong class is
    winning), the probability is high---urgent correction. When
    $v_k$ is low or negative, the probability drops.
\end{itemize}

\begin{keyinsight}
The voting margin directly controls the training intensity. Clauses
that are already doing well receive little feedback; clauses that are
struggling receive intense feedback. This is an automatic form of
\emph{hard example mining}---the system focuses its learning effort
where it is most needed.
\end{keyinsight}


% ----------------------------------------------------------------------------
\section{Worked Example: Learning a Simple Pattern}
\label{sec:worked-example}
% ----------------------------------------------------------------------------

Let us trace the training of a single positive clause for class~``1''
on a tiny 4-pixel image. We use $N = 3$ states per action (6 states
total per automaton) and $s = 3.0$.

\subsection{Setup}

The input has 4 pixels, giving 8 literals after expansion:
$x_1, \bar{x}_1, x_2, \bar{x}_2, x_3, \bar{x}_3, x_4, \bar{x}_4$.

We assign one automaton to each literal, all starting in state~4
(just above the boundary, action = include). So initially the clause
includes all 8 literals: $C = x_1 \wedge \bar{x}_1 \wedge \cdots$,
which always outputs 0 (since $x_k \wedge \bar{x}_k = 0$).

\begin{center}
\begin{tabular}{@{}l*{8}{c}@{}}
  \toprule
  \textbf{Literal} & $x_1$ & $\bar{x}_1$ & $x_2$ & $\bar{x}_2$ &
    $x_3$ & $\bar{x}_3$ & $x_4$ & $\bar{x}_4$ \\
  \midrule
  Initial state & 4 & 4 & 4 & 4 & 4 & 4 & 4 & 4 \\
  Action & inc & inc & inc & inc & inc & inc & inc & inc \\
  \bottomrule
\end{tabular}
\end{center}

\subsection{Training Example 1}

Input: image of ``1'' with pixels $(1, 0, 1, 0)$, so:
\[
  (x_1, \bar{x}_1, x_2, \bar{x}_2, x_3, \bar{x}_3, x_4, \bar{x}_4)
  = (1, 0, 0, 1, 1, 0, 0, 1)
\]

The clause output is 0 (since many included literals are 0). Since
$C = 0$ and this is a positive clause for the true class, Type~I
feedback applies with the ``clause output = 0'' rules.

Each automaton receives a penalty with probability $\frac{1}{s} =
\frac{1}{3} \approx 0.33$. Suppose random draws penalise automata
for $\bar{x}_1$, $x_2$, $\bar{x}_3$, and $x_4$. These automata
move from state~4 to state~3 (crossing the boundary from include to
exclude):

\begin{center}
\begin{tabular}{@{}l*{8}{c}@{}}
  \toprule
  \textbf{Literal} & $x_1$ & $\bar{x}_1$ & $x_2$ & $\bar{x}_2$ &
    $x_3$ & $\bar{x}_3$ & $x_4$ & $\bar{x}_4$ \\
  \midrule
  State & 4 & \textbf{3} & \textbf{3} & 4 & 4 & \textbf{3} & \textbf{3} & 4 \\
  Action & inc & \textbf{exc} & \textbf{exc} & inc & inc & \textbf{exc} & \textbf{exc} & inc \\
  \bottomrule
\end{tabular}
\end{center}

Now the clause is $C = x_1 \wedge \bar{x}_2 \wedge x_3 \wedge
\bar{x}_4$. Let us check: for the input $(1, 0, 1, 0)$, the literal
values at included positions are $x_1 = 1$, $\bar{x}_2 = 1$,
$x_3 = 1$, $\bar{x}_4 = 1$. All are 1, so $C = 1$. The clause
now matches this ``1'' pattern.

\subsection{Training Example 2}

Input: image of ``0'' with pixels $(1, 1, 1, 1)$, so all positive
literals are 1 and all negated literals are 0. The clause evaluates:
$x_1 = 1$, $\bar{x}_2 = 0$---stop, $C = 0$. Good: the clause does
not match this ``0'' image.

But now suppose another ``0'' image arrives with pixels $(1, 0, 1, 1)$:
$x_1 = 1$, $\bar{x}_2 = 1$, $x_3 = 1$, $\bar{x}_4 = 0$---stop,
$C = 0$. Still correct.

And a ``1'' image with pixels $(1, 0, 0, 0)$:
$x_1 = 1$, $\bar{x}_2 = 1$, $x_3 = 0$---stop, $C = 0$. This ``1''
pattern does not match because $x_3 = 0$. This clause has learned a
\emph{specific} pattern for ``1'' that requires both $x_1$ and $x_3$
to be bright.

\subsection{Continuing Training}

Over many examples, Type~I feedback will:
\begin{itemize}
  \item Strongly reinforce included literals that are consistently
    true for class ``1'' images (pushing automata toward state~6).
  \item Gently push excluded literals further into exclusion (toward
    state~1).
\end{itemize}

If a Type~II signal arrives (because this clause wrongly matches an
image from another class), it will try to include a literal that is
0 in that image, preventing the false match.


% ----------------------------------------------------------------------------
\section{The Role of Multiple Clauses}
\label{sec:multiple-clauses}
% ----------------------------------------------------------------------------

A single clause can only detect one specific pattern. The digit ``3''
can be written in many ways---some with a flat top, some with a
rounded top, some thick, some thin. No single conjunctive clause can
match all variants.

This is why each class has $m$ clauses (typically $m = 100$ to
$2000$). Each clause specialises in a different sub-pattern:

\begin{itemize}
  \item Clause~1 might detect ``3'' with a flat top bar.
  \item Clause~2 might detect ``3'' with a rounded top.
  \item Clause~3 might detect a common stroke in the middle.
  \item And so on.
\end{itemize}

The vote aggregation (sum of clause outputs) acts as a soft OR: if
\emph{any} clause matches, the vote for class ``3'' increases.

\begin{analogy}
Imagine a committee of experts, each looking at a different aspect of
the digit. Expert~1 checks if the top stroke looks right. Expert~2
checks the middle junction. Expert~3 checks the bottom curve. They
vote, and the digit is classified based on the majority opinion. The
Tsetlin Machine learns both the experts' criteria (which pixels to
check) and which experts to trust (positive vs.\ negative polarity).
\end{analogy}


% ----------------------------------------------------------------------------
\section{Negative Clauses: Learning What a Digit Is Not}
\label{sec:negative-clauses}
% ----------------------------------------------------------------------------

Negative clauses vote \emph{against} their class. A negative clause
for class ``3'' might detect patterns that look like ``8'' (which is
similar to ``3'' but has a closed top loop). When this pattern is
present, the negative clause fires and reduces the vote for ``3,''
making it less likely to be confused with ``8.''

\begin{example}
Suppose class ``7'' has:
\begin{itemize}
  \item A positive clause that fires when there is a horizontal top
    stroke and a diagonal down stroke.
  \item A negative clause that fires when there is also a horizontal
    middle stroke (suggesting ``1'' written with a crossbar, or
    some styles of ``4'').
\end{itemize}
The positive clause says ``this might be a 7'' and the negative clause
says ``but this feature argues against 7.'' The vote is the net
effect.
\end{example}


% ----------------------------------------------------------------------------
\section{Hyperparameters and Their Effects}
\label{sec:hyperparams}
% ----------------------------------------------------------------------------

The Tsetlin Machine has remarkably few hyperparameters:

\begin{table}[htbp]
\centering
\caption{Tsetlin Machine hyperparameters.}
\label{tab:tm-hyperparams}
\begin{tabular}{@{}llp{6cm}@{}}
  \toprule
  \textbf{Symbol} & \textbf{Typical range} & \textbf{Effect} \\
  \midrule
  $m$ & 100--2000 & Number of clauses per class. More clauses
    $\Rightarrow$ more capacity to learn diverse patterns, but
    slower training and inference. \\
  $T$ & 10--100 & Vote threshold. Controls how many clauses
    actively train. Lower $T$ $\Rightarrow$ fewer active clauses,
    faster convergence but potentially lower accuracy. \\
  $s$ & 2.0--10.0 & Specificity. Higher $s$ $\Rightarrow$ more
    specific clauses (more included literals). Lower $s$ $\Rightarrow$
    more general clauses. \\
  $N$ & 50--200 & States per action in each automaton. Higher $N$
    $\Rightarrow$ more noise resistance, slower adaptation. \\
  \bottomrule
\end{tabular}
\end{table}

\begin{engineernote}
A good starting configuration for MNIST with simple thresholding:
$m = 500$ clauses per class (250 positive, 250 negative), $T = 50$,
$s = 5.0$, $N = 100$. This achieves approximately 96--97\% test
accuracy and trains in under a minute on a modern CPU.
\end{engineernote}


% ----------------------------------------------------------------------------
\section{Why Stochastic Feedback Creates Diversity}
\label{sec:diversity}
% ----------------------------------------------------------------------------

A crucial question: if all clauses for class ``3'' receive the same
training examples, why do they learn different patterns?

The answer is the stochastic feedback. Each automaton independently
receives or does not receive a reward/penalty based on random number
generation. This means:

\begin{enumerate}
  \item Two clauses seeing the same input receive different feedback
    on different literals.
  \item Over time, random variations cause clauses to diverge: one
    clause might include $x_{42}$ while another excludes it, leading
    them to specialise on different sub-patterns.
  \item The vote threshold $T$ amplifies this effect: once enough
    clauses are matching, others stop receiving feedback and are free
    to explore alternative patterns.
\end{enumerate}

\begin{keyinsight}
Diversity in the Tsetlin Machine is not engineered---it is an emergent
property of stochastic feedback combined with vote thresholding. Each
clause follows its own random walk through literal space, and the
aggregate system self-organises into a diverse committee of pattern
detectors.
\end{keyinsight}


% ----------------------------------------------------------------------------
\section{Comparison with Gradient-Based Learning}
\label{sec:tm-vs-gradient}
% ----------------------------------------------------------------------------

It is instructive to compare the Tsetlin Machine training loop with
standard gradient descent:

\begin{table}[htbp]
\centering
\caption{Tsetlin Machine vs.\ gradient descent.}
\label{tab:tm-vs-gd}
\begin{tabular}{@{}lll@{}}
  \toprule
  \textbf{Aspect} & \textbf{Gradient descent} & \textbf{Tsetlin Machine} \\
  \midrule
  Parameter type & Floating-point weights & Integer automaton states \\
  Update rule & $w \leftarrow w - \eta \nabla$ & $s \leftarrow s \pm 1$ \\
  Requires & Differentiable ops & Any reward/penalty signal \\
  Arithmetic & FP multiply-add & Integer increment/decrement \\
  Hyperparams & Learning rate, optimiser, & $s$, $T$, $m$, $N$ \\
   & momentum, weight decay & \\
  Memory & 32 bits per weight & $\lceil \log_2(2N) \rceil$ bits per automaton \\
  Interpretability & Opaque weight matrices & Readable Boolean clauses \\
  \bottomrule
\end{tabular}
\end{table}

The Tsetlin Machine's main advantages are simplicity, interpretability,
and integer-only arithmetic~\cite{wheeldon2020pervasive}. Its main disadvantage is slower convergence
on complex tasks~\cite{zhang2021convergence}, which is partly why we combine it with a Logic Gate
Network~\cite{petersen2022difflogic} in the hybrid architecture.


% ----------------------------------------------------------------------------
\section{Summary}
\label{sec:ch6-summary}
% ----------------------------------------------------------------------------

The Tsetlin Machine training algorithm can be summarised in four
principles:

\begin{enumerate}
  \item \textbf{Type~I reinforces:} For the correct class, strengthen
    clauses that match the input by rewarding included-and-true
    literals and penalising excluded-but-true ones.
  \item \textbf{Type~II sabotages:} For incorrect classes, break
    wrongly-matching clauses by including literals that are false
    in the current input.
  \item \textbf{Stochastic diversity:} Random feedback probabilities
    ensure different clauses learn different patterns.
  \item \textbf{Vote thresholding:} The threshold $T$ balances
    training effort across clauses, preventing any single pattern
    from monopolising the vote.
\end{enumerate}

The next chapter applies these principles to MNIST and walks through
the training process on real digit images.
